\section{Conclusions}\label{sec:conclusions}
\paragraph{How the tasks relate.} The tasks form a chain. T2's lengths are the weights of every shortest path; T3 turns them into sectors; T4 reads the oldest-water node of each sector from the same Dijkstra tree; T5 uses that node as the sink of a flow problem restricted to the sector. T6 reuses T2's all-pairs distances as its metric closure, and T1 decides the representation all of them run on. The results also explain each other. The leaf sources found in the data stage are both the single-pipe bottlenecks of T5 and the dead ends the T6 walk must traverse twice. PES's extreme circuity in T2 reflects two clusters joined by few pipes, which is why separating its sectors needs only two closures in T3. T4 and T5 on the same source--sink pairs show concretely that the shortest route and the maximum flow are different problems: the route minimises length, the flow is set by the narrowest cut, and in the meshed FOS the flow is about three times what the shortest route alone could carry.

\paragraph{Cost and difficulty.} T6 is the only hard problem. Its exact solution grows exponentially with the number of locations, so optimality was proven only for NYT and for subsets of at most 24 nodes; for HAN, FOS and PES the routes are heuristic. Among the polynomial steps, Floyd-Warshall ($O(V^3)$) is the most expensive, and the adjacency matrix in general costs $O(V^2)$ where the list costs $O(V+E)$. Dijkstra, the flow algorithms and the nearest-node search are all fast at this scale; the flow algorithms run far below their worst-case bounds because the sectors need at most four augmenting paths. The cheap heuristics carry most of the value of the expensive method: 2-opt was within 1.8\% of the optimum on average. A selection hyper-heuristic over the same kind of moves made the heuristic routes more reliable across runs and shortened the best PES route by 2.3\%, still without a guarantee.

\paragraph{Limitations.} The model is purely graph-theoretic: there is no hydraulics (pressure, head loss, demand), capacities are taken as given numbers without units, and water age is taken as proportional to distance rather than to travel time. Flow is analysed for one source--sink pair per sector inside the closed sector, not for the whole network. Several conventions were needed and are stated where used: the treatment of PES's repeated pipe, the rejection of NYT's outlying new point, the sequential insertion of new nodes, and tie-breaking towards the earlier source (no ties occurred). The branch-and-bound bound is simple and weak on near-tree networks, so the full HAN, FOS and PES routes remain uncertified. Finally, the instances are small ($V\le71$), so performance conclusions rest on growth across subsets and are affected by Python interpreter overhead.
